\documentclass[11pt]{article}

\usepackage[T1]{fontenc}
\usepackage{lmodern}
\usepackage{microtype}
\usepackage{amsmath,amssymb,amsthm}
\usepackage{booktabs}
\usepackage{tabularx}
\usepackage{array}
\usepackage{enumitem}
\usepackage[table]{xcolor}
\usepackage[margin=1in]{geometry}
\usepackage{xurl}
\usepackage[
  colorlinks=true,
  linkcolor=blue!55!black,
  citecolor=blue!55!black,
  urlcolor=blue!55!black,
  pdfauthor={Austin Lin Gibbons},
  pdftitle={An Order-94 Rowley Template and the Bound R5(5) >= 42078}
]{hyperref}
\usepackage{fancyhdr}

\definecolor{headingblue}{RGB}{31,77,120}
\definecolor{tableblue}{RGB}{38,82,122}
\definecolor{lightrow}{RGB}{242,246,250}

\hypersetup{bookmarksopen=true}
\setlist[itemize]{leftmargin=1.5em,itemsep=2pt,topsep=4pt}
\setlist[enumerate]{leftmargin=1.7em,itemsep=2pt,topsep=4pt}
\setlength{\parindent}{0pt}
\setlength{\parskip}{5pt plus 1pt minus 1pt}
\renewcommand{\arraystretch}{1.12}
\clubpenalty=10000
\widowpenalty=10000
\displaywidowpenalty=10000

\makeatletter
\renewcommand{\@seccntformat}[1]{%
  \csname the#1\endcsname\hspace{0.8em}}
\makeatother

\pagestyle{fancy}
\fancyhf{}
\fancyhead[L]{\small An order-94 Rowley template}
\fancyhead[R]{\small\thepage}
\renewcommand{\headrulewidth}{0.3pt}

\newtheorem{theorem}{Theorem}[section]
\newtheorem{proposition}[theorem]{Proposition}
\newtheorem{lemma}[theorem]{Lemma}
\newtheorem{corollary}[theorem]{Corollary}
\theoremstyle{definition}
\newtheorem{definition}[theorem]{Definition}
\newtheorem{certificate}[theorem]{Certificate}
\theoremstyle{remark}
\newtheorem{remark}[theorem]{Remark}

\newcommand{\R}{\mathbb{R}}
\newcommand{\Z}{\mathbb{Z}}
\newcommand{\K}{K}
\newcommand{\Rfive}{R_5(5)}
\newcommand{\code}[1]{\texttt{#1}}
\newcommand{\sha}[1]{\begingroup\ttfamily\scriptsize\urlstyle{same}\nolinkurl{#1}\endgroup}
\newcommand{\hashbreak}[2]{{\ttfamily\scriptsize #1\linebreak #2}}

\title{\bfseries An Order-94 Rowley Template and the Bound\\[3pt]
  \texorpdfstring{\boldmath$\Rfive\ge 42078$}{R5(5) >= 42078}}
\author{Austin Lin Gibbons\thanks{Draft for independent review. Mathematical
correctness, priority, and publication readiness are treated as separate
questions throughout.}}
\date{23 July 2026}

\begin{document}
\maketitle

\begin{abstract}
We give an explicit effective Rowley $(5,5,3)$ template of order $94$ with
extension parameter $\phi=40$.  Its two inherited distance colours remain
$K_5$-free under period-$93$ repetition through the exact compression cutoff
$4(94-2)=368$, while its third distance class is a sum-free triangle-free
template.  We combine this template with an explicit cyclic three-colouring
of $K_{453}$ recovered from Rowley's archived ancillary data.  The resulting
five distance classes partition $\{1,\ldots,42076\}$ and define a cyclic
five-colouring of $K_{42077}$.  A self-contained specialization of Rowley's
composition argument proves that none of the five colours contains $K_5$.
Consequently,
\[
  R_5(5)\ge 42078.
\]
All finite claims are checked by exact algorithms.  The order-$94$ template
and order-$453$ prototype are each accepted by independent Python and C++
clique searches using different internal representations, and the full
compound word is reconstructed independently from the displayed set-union
formula.  The checked revision of the standard survey records the previous
bound $R_5(5)\ge41626$; priority beyond the literature corpus described here
remains subject to external review.
\end{abstract}

\medskip
\noindent\textbf{Keywords.} Multicolour Ramsey numbers; linear colourings;
circulant colourings; triangle-free templates; exact computation.

\noindent\textbf{MSC 2020.} 05C55, 05D10, 68W30.

\section{Introduction}

For positive integers $r$ and $k$, the diagonal multicolour Ramsey number
$R_r(k)$ is the least $N$ such that every $r$-colouring of the edges of
$K_N$ contains a monochromatic $K_k$.  An explicit $r$-colouring of
$K_{N-1}$ without monochromatic $K_k$ proves $R_r(k)\ge N$.

Rowley developed a composition method in which a small linear colouring
containing a distinguished triangle-free template is combined with a second
linear Ramsey colouring \cite{rowley-general,rowley-sat}.  In the case
relevant here, his published order-$93$ effective $(5,5,3)$ template, with
$\phi=40$, combines with an order-$453$ three-colouring to give a
five-colouring of order $41625$, and hence
\[
  R_5(5)\ge41626.
\]
This value is recorded in revision 18 of Radziszowski's survey
\cite{radziszowski}.

The contribution of this paper is one additional template position.  We
exhibit an order-$94$ effective $(5,5,3)$ template with the same value
$\phi=40$.  The compound order becomes
\[
  (94-1)(453-1)+1+40
  =93\cdot452+41
  =42077.
\]
Thus the one-position extension is amplified into a gain of $452$ vertices.

The proof has three logically separate inputs.
\begin{itemize}
  \item Section~\ref{sec:template} gives the explicit order-$94$ word and
  the exact finite conditions it satisfies.
  \item Section~\ref{sec:prototype} gives the explicit order-$453$
  prototype recovered from Rowley's archived 2022 ancillary data.
  \item Sections~\ref{sec:compound} and~\ref{sec:proof} define the compound
  and prove directly that its five colours are $K_5$-free.
\end{itemize}
Search heuristics are not part of the proof.  The frozen words, the
composition formula, and exact checkers form the certificate.

\begin{remark}[Status and scope]
The theorem proved here is a correctness claim.  The specific value $42078$
was not found in the checked literature, and the checked survey revision
records $41626$, but those facts do not establish priority.  A public
priority claim should await an independent literature search and contact
with specialists familiar with unpublished linear Ramsey constructions.
\end{remark}

\section{Linear colourings and triangle-free templates}
\label{sec:definitions}

\begin{definition}[Linear distance colouring]
For $m\ge2$, a linear $q$-colouring of $K_m$ is a map
\[
  c:\{1,\ldots,m-1\}\longrightarrow\{1,\ldots,q\}.
\]
The edge $\{i,j\}$, where $0\le i<j\le m-1$, receives colour $c(j-i)$.
It is cyclic if $c(d)=c(m-d)$ for every $1\le d\le m-1$.
\end{definition}

For a three-colour word $c$, put
\[
  A_1=c^{-1}(1),\qquad A_2=c^{-1}(2),\qquad T=c^{-1}(3).
\]

\begin{definition}[Triangle-free template]
The third distance class $T$ is a triangle-free template, or
\emph{tf-template}, if $m-1\in T$ and the colour-$3$ graph is triangle-free.
Its extension parameter is
\[
  \phi=\min(T)-1.
\]
\end{definition}

\begin{lemma}[Sum-free criterion]
\label{lem:sumfree}
The colour-$3$ graph of a linear word is triangle-free if and only if there
do not exist positive integers $a,b$ such that
\[
  a\in T,\qquad b\in T,\qquad a+b\in T.
\]
The case $a=b$ is included.
\end{lemma}

\begin{proof}
Translate the least vertex of any triangle to $0$ and write the other two
vertices as $a<a+b$.  Its three edge lengths are $a$, $b$, and $a+b$.
The converse is immediate from the same three vertices.
\end{proof}

For a period $p=m-1$, the inherited colours are repeated by assigning a
positive distance $d$ the base colour
\begin{equation}
  c\!\left(((d-1)\bmod p)+1\right).
  \label{eq:periodic-colour}
\end{equation}
Thus residue zero is represented by base distance $p$, not by $0$.  This
convention is essential because the terminal distance $p$ belongs to the
template colour.

\begin{definition}[Effective $(5,5,3)$ certificate]
\label{def:effective}
An order-$m$ word is an effective $(5,5,3)$ certificate with parameter
$\phi$ if:
\begin{enumerate}
  \item $m-1\in T$ and $\min(T)-1=\phi$;
  \item $T$ satisfies Lemma~\ref{lem:sumfree}; and
  \item for each $s\in\{1,2\}$, the periodic colour defined by
  \eqref{eq:periodic-colour} has no $K_5$ on vertices
  \[
    0,1,\ldots,4(m-2).
  \]
\end{enumerate}
\end{definition}

The final cutoff is the specialization to $K_5$ of the compression argument
in Rowley's Generalised Construction Theorem
\cite[Definition 3.1 and Theorem 3.2]{rowley-general}.  For completeness,
the same compression is written explicitly in
Section~\ref{sec:proof}.

\section{The order-94 template}
\label{sec:template}

\begin{certificate}[Canonical order-$94$ word]
\label{cert:order94}
Let $c(1)c(2)\cdots c(93)$ be the following word; displayed line breaks are
typographic only:
\begin{quote}
\ttfamily\footnotesize
11121221222212111212212111212222122121113131232132323333333333333333333333\\
3333333232312321313
\end{quote}
Its distance classes are
\begin{align*}
A_1={}&\{1,2,3,5,8,13,15,16,17,19,22,24,25,26,28,33,36,38,39,40,\\
&\hspace{1.75em}42,44,48,86,90,92\},\\
A_2={}&\{4,6,7,9,10,11,12,14,18,20,21,23,27,29,30,31,32,34,35,37,\\
&\hspace{1.75em}45,47,50,52,82,84,87,89\},\\
T={}&\{41,43,46,49,51,53,54,55,56,57,58,59,60,61,62,63,64,65,66,\\
&\hspace{1.75em}67,68,69,70,71,72,73,74,75,76,77,78,79,80,81,83,85,88,91,93\}.
\end{align*}
\end{certificate}

\begin{proposition}[Exact finite properties]
\label{prop:order94}
Certificate~\ref{cert:order94} is an effective $(5,5,3)$ certificate of
order $94$ with $\phi=40$.  More explicitly:
\begin{enumerate}
  \item $93\in T$, $T\cap\{1,\ldots,40\}=\varnothing$, and $\min(T)=41$;
  \item no $a,b,a+b$ all lie in $T$, including the case $a=b$;
  \item neither inherited colour contains $K_5$ on vertices
  $0,\ldots,368$; and
  \item the reflection identities
  \[
    c(d)=c(41-d)\quad(1\le d\le40),
    \qquad
    c(d)=c(134-d)\quad(41\le d\le93)
  \]
  hold.
\end{enumerate}
\end{proposition}

\begin{proof}
Items (1), (2), and (4) are direct complete enumerations on the displayed
word.  For item (3), two independent exact clique searches construct the
periodic colour graphs on the $369$ vertices $0,\ldots,368$ and exhaustively
reject $K_5$ in both inherited colours.  The Python checker uses arbitrary
precision integer bit masks for forward-neighbour sets; the C++ checker uses
ordered integer candidate vectors and pairwise filtering.  Their algorithms
and regression coverage are described in Section~\ref{sec:verification}.
\end{proof}

The canonical file is \code{results/order94\_t12.template}, with SHA-256
\begin{center}
\sha{a3dd3415956277d68fc0197f6f5e35c38f2ed7936c66f11f658b0c57688abec1}.
\end{center}
The same word is also accepted when overtested through distance $372$.
Two further distinct accepted order-$94$ words are retained in the artifact
package as robustness evidence, but they are not used below.

\section{The order-453 prototype}
\label{sec:prototype}

Rowley's 2022 construction starts from a cyclic order-$n$ prototype, copies
its colours at distances $1,\ldots,n-1$, and places the new template colour
at distance $n$ \cite[Section 5]{rowley-sat}.  The archived order-$977$
$(5,5,5,3)$ template in ancillary sheet \code{Paper\_Sep\_2022}, column
\code{AB}, was obtained by extending an Exoo-derived order-$453$
$(5,5,5)$ graph.  Consequently, entries $1,\ldots,452$ of that column are
the prototype word, and entry $453$ begins colour $4$.

The full $452$-symbol word is printed in Appendix~\ref{app:words}.

\begin{proposition}[Source-recovered prototype]
\label{prop:prototype}
The word in Appendix~\ref{app:words} defines a cyclic three-colouring of
$K_{453}$ with no monochromatic $K_5$.
\end{proposition}

\begin{proof}
Every symbol was compared with the SHA-256-fixed ancillary XML.  The
comparison also checks the sheet metadata and confirms that distance $453$
begins the new template colour.  Independent Python bit-mask and C++
ordered-vector searches then exhaustively reject $K_5$ in each of the three
finite colour graphs.  Finally, direct comparison gives
$v(d)=v(453-d)$ for all $1\le d\le452$.
\end{proof}

The extracted artifact is
\code{sources/rowley\_exoo\_order453.prototype}, with SHA-256
\begin{center}
\sha{19c97e6279c184f6f462786cadda4b7c9773d870a5b680a04eb1503ef384a2d0}.
\end{center}
Once the word is explicit and Proposition~\ref{prop:prototype} is verified,
the historical attribution is not a logical dependency of the main theorem.

\section{The compound colouring}
\label{sec:compound}

Set $p=93$.  Let $A_1,A_2,T$ be the distance classes in
Certificate~\ref{cert:order94}, and let $B_1,B_2,B_3$ be the three distance
classes of the order-$453$ prototype.

For $s\in\{1,2\}$, define the inherited output classes
\begin{equation}
\begin{split}
A'_s={}&
\{\ell+(\mu-1)p:\ell\in A_s,\ 1\le\mu\le452\}\\
&{}\cup
\{\ell+452p:\ell\in A_s,\ 1\le\ell\le40\}.
\end{split}
\label{eq:inherited-sets}
\end{equation}
For $j\in\{1,2,3\}$, relabel the three prototype colours as output colours
$j+2$ and put
\begin{equation}
  B'_j=\{\ell+(\mu-1)p:\ell\in T,\ \mu\in B_j\}.
  \label{eq:prototype-sets}
\end{equation}

\begin{lemma}[Partition]
\label{lem:partition}
The five sets $A'_1,A'_2,B'_1,B'_2,B'_3$ partition
$\{1,\ldots,42076\}$.
\end{lemma}

\begin{proof}
First consider $1\le d\le452p$.  Write $d$ uniquely as
$d=\ell+(\mu-1)p$ with $1\le\ell\le p$ and $1\le\mu\le452$.
If $\ell\in A_s$, then $d\in A'_s$.  If $\ell\in T$, then $\mu$ lies in
exactly one prototype class $B_j$, and $d\in B'_j$.  These cases are
disjoint and exhaustive.

The remaining distances are $452p+\ell$ with $1\le\ell\le40$.  Since
$T\cap\{1,\ldots,40\}=\varnothing$, each such $\ell$ lies in exactly one
of $A_1,A_2$, and the second union in \eqref{eq:inherited-sets} assigns the
distance uniquely.
\end{proof}

Colour the edge $\{x,y\}$, where $0\le x<y\le42076$, by the unique one of
the five sets in Lemma~\ref{lem:partition} containing $y-x$.  This is the
complete explicit colouring.

\begin{table}[htbp]
\centering
\caption{Distance-class sizes in the compound colouring.}
\label{tab:class-sizes}
\begin{tabular}{@{}ccc@{}}
\toprule
Output colour & Number of distances & Origin\\
\midrule
1 & 11772 & first inherited template colour\\
2 & 12676 & second inherited template colour\\
3 & 5304  & prototype colour 1\\
4 & 7566  & prototype colour 2\\
5 & 4758  & prototype colour 3\\
\midrule
Total & 42076 & all positive distances\\
\bottomrule
\end{tabular}
\end{table}

\section{Proof of the Ramsey bound}
\label{sec:proof}

\begin{theorem}
\label{thm:main}
There is a five-colouring of $K_{42077}$ with no monochromatic $K_5$.
Consequently,
\[
  R_5(5)\ge42078.
\]
\end{theorem}

\begin{proof}
We separately exclude cliques in inherited colours and in
prototype-derived colours.

\emph{Inherited colours.}
Fix $s\in\{1,2\}$ and suppose that a colour-$s$ $K_5$ exists in the
compound.  Translate its least vertex to $0$ and write the other vertices
as
\[
  0<h_1<h_2<h_3<h_4.
\]
Every $h_i$ and every difference $h_j-h_i$ has a residue in $A_s$ modulo
$p=93$.

If $h_1>p$, subtract $p$ from all four positive coordinates.  If a
consecutive gap $h_{i+1}-h_i$ exceeds $p$, subtract $p$ from the entire
upper tail $h_{i+1},\ldots,h_4$.  These operations preserve positivity,
strict order, and every relevant colour: differences across the cut change
by exactly $p$, while differences within either part do not change.

Iteration terminates with $h_1\le p$ and every consecutive gap at most
$p$.  No inherited-colour edge can have length divisible by $p$, because
the representative of residue zero is $p=93\in T$.  Hence
$h_1\le p-1$, every consecutive gap is at most $p-1$, and
\[
  h_4\le4(p-1)=368.
\]
This gives a $K_5$ in the exact periodic graph rejected by
Proposition~\ref{prop:order94}, a contradiction.

\emph{Prototype-derived colours.}
Fix $j\in\{1,2,3\}$ and suppose that output colour $j+2$ contains a
$K_5$.  Translate its least vertex to $0$.  Each other vertex has a unique
form
\[
  d_i=t_i+(\mu_i-1)p,
  \qquad t_i\in T,\quad \mu_i\in B_j.
\]
Take two vertices with $d_i<d_k$.  If $t_i<t_k$, then the positive residue
of $d_k-d_i$ is $t_k-t_i$.  This residue cannot lie in $T$, because
\[
  t_i+(t_k-t_i)=t_k
\]
would contradict the sum-free property of $T$.  Thus an edge in the same
prototype-derived colour forces $t_i\ge t_k$.

If $t_i=t_k$, the difference has prototype index $\mu_k-\mu_i$.  If
$t_i>t_k$, then
\[
  d_k-d_i=
  \bigl(p-(t_i-t_k)\bigr)+(\mu_k-\mu_i-1)p,
\]
whose prototype index is again $\mu_k-\mu_i$.  Membership of every
pairwise difference in output colour $j+2$ therefore implies
\[
  \mu_k-\mu_i\in B_j.
\]
The indices $\mu_i$ are strictly increasing because the $d_i$ increase
while the $t_i$ do not.  Hence
\[
  \{0,\mu_1,\mu_2,\mu_3,\mu_4\}
\]
is a monochromatic $K_5$ in prototype colour $j$, contradicting
Proposition~\ref{prop:prototype}.

All five output colours are therefore $K_5$-free.  The compound has order
\[
  93\cdot452+41=42077,
\]
so the definition of $R_5(5)$ gives $R_5(5)\ge42078$.
\end{proof}

\begin{corollary}
The linear colouring in Theorem~\ref{thm:main} is cyclic.
\end{corollary}

\begin{proof}
The prototype is cyclic, and Certificate~\ref{cert:order94} satisfies
Rowley's two reflection identities stated in
Proposition~\ref{prop:order94}(4).  Direct comparison of the expanded word
also gives $w(d)=w(42077-d)$ for every $1\le d\le42076$.
\end{proof}

\section{Exact computational verification}
\label{sec:verification}

The computational claims in Propositions~\ref{prop:order94} and
\ref{prop:prototype} are finite, exhaustive, and deterministic.  No random
search decision enters their verdicts.

\subsection{Clique search}

Both implementations use the following exhaustive recursion, with different
representations of the candidate set.

\begin{quote}
\ttfamily\small
Search(prefix, candidates, need):\\
\phantom{xx}if need = 0: return prefix\\
\phantom{xx}if size(candidates) < need: return NONE\\
\phantom{xx}while size(candidates) >= need:\\
\phantom{xxxx}v = least vertex in candidates\\
\phantom{xxxx}delete v from candidates\\
\phantom{xxxx}next = \{u in candidates : Adj(v,u)\}\\
\phantom{xxxx}result = Search(prefix followed by v, next, need-1)\\
\phantom{xxxx}if result is not NONE: return result\\
\phantom{xx}return NONE
\end{quote}

Every candidate list is strictly increasing.  At each call it consists
exactly of the larger vertices adjacent to every prefix vertex.  Thus every
clique is represented once.  The cardinality test prunes only branches that
cannot reach the target size, so returning \code{NONE} is an exhaustive
proof of nonexistence for the supplied finite graph.

\subsection{Independent implementations and finite counts}

The Python implementation stores forward-neighbour sets in arbitrary
precision integer bit masks.  The C++ implementation stores ordered integer
vectors and forms intersections by explicit pairwise filtering.  They share
the written mathematical semantics but neither parser nor clique-search
code.

\begin{table}[htbp]
\centering
\caption{Exact search-node counts.  Different counts reflect different
candidate orderings and data structures.}
\label{tab:nodes}
\begin{tabular}{@{}lrrr@{}}
\toprule
Finite object and checker & Colour 1 & Colour 2 & Colour 3\\
\midrule
Order-$94$ template, Python & 550261 & 718490 & --\\
Order-$453$ prototype, Python & 1040265 & 5811421 & 782585\\
Order-$453$ prototype, C++ & 237855 & 837235 & 201199\\
\bottomrule
\end{tabular}
\end{table}

Seven negative fixtures test malformed length, a missing terminal template
distance, forbidden-prefix use, an additive template triple, insufficient
repeat span, and explicit repeated $K_5$ witnesses in each inherited colour.
Both template implementations make the same accept/reject decisions on the
complete fixture suite.

\subsection{Source extraction and compound reconstruction}

The order-$93$ published seed and the order-$453$ prototype are compared
symbol for symbol with the ancillary XML in the fixed arXiv source archive.
The archive SHA-256 is
\begin{center}
\sha{5041096a33d6898c116a35a102e521fb14fe872fb18c20ff7d30822ed4c396b2}.
\end{center}

The final expanded word is generated once by a block/residue implementation
of \eqref{eq:inherited-sets} and \eqref{eq:prototype-sets}.  A second program,
which does not import the generator, constructs the five displayed sets as
literal unions, verifies that every distance $1,\ldots,42076$ occurs exactly
once, compares every symbol with the frozen word, checks cyclic reflection,
and confirms SHA-256
\begin{center}
\sha{a97d5dce8a927db2889ba220119f9d4f3d1b88ee24d441f82a803a195ae8028d}.
\end{center}

\begin{remark}
The $42077$-vertex graph was not subjected to a naive all-vertex clique
search.  Its $K_5$-freeness is established by the two preservation arguments
in Theorem~\ref{thm:main}, whose finite hypotheses are checked exactly.
\end{remark}

\section{Reproducibility and review protocol}
\label{sec:repro}

The complete deterministic verification is run from the artifact package by
\begin{quote}
\ttfamily\small
sh tests/run\_end\_to\_end\_checks.sh
\end{quote}
The suite:
\begin{enumerate}
  \item re-extracts the published order-$93$ seed and order-$453$ prototype
  from the fixed ancillary XML;
  \item verifies all three prototype colours with the Python and C++
  clique searches;
  \item verifies three distinct order-$94$ words with the Python and C++
  full-condition checkers;
  \item overtests the canonical word through distance $372$; and
  \item reconstructs all $42076$ compound distances independently and
  checks the frozen hash.
\end{enumerate}

An adversarial reproduction should pay particular attention to:
\begin{itemize}
  \item the representative $93$ for residue zero;
  \item inclusion of $a=b$ in the sum-free test;
  \item the distinction between maximum distance $368$ and vertex count
  $369$;
  \item the source-theorem wording caveat concerning the aggregate cutoff,
  versus the explicit per-colour bound in its proof; and
  \item the recovery of distances $1,\ldots,452$ from ancillary column
  \code{AB}, with distance $453$ beginning template colour $4$.
\end{itemize}

The remaining nonmathematical gates are an independent-machine rerun, a
clean-room third checker if desired, and a specialist priority search.

\section{Literature status}

Rowley's 2022 paper records the order-$93$ template and the resulting bound
$R_5(5)\ge41626$ \cite{rowley-sat}.  Revision 18 of the standard dynamic
survey, dated 24 April 2026, records the same value in Table XIa
\cite{radziszowski}.  The January 2026 asymptotic update of Campos and
Pohoata does not state this specific improvement \cite{campos-pohoata}.

Searches for the value $42078$, an order-$94$ $(5,5,3)$ Rowley template,
and follow-up work to the cited construction papers found no collision in
the checked corpus.  The appropriate status is therefore:
\begin{center}
\textbf{not found in the checked literature; external priority review
required.}
\end{center}

\appendix

\section{Complete certificate words}
\label{app:words}

\subsection{Order-94 template}

For convenience, the canonical order-$94$ word is repeated here:
\begin{quote}
\ttfamily\footnotesize
11121221222212111212212111212222122121113131232132323333333333333333333333\\
3333333232312321313
\end{quote}

\subsection{Order-453 prototype}

The following seven lines are to be concatenated exactly:
\begin{quote}
\ttfamily\scriptsize
12223333221211121312122313122213232232332221323223233221111122121112213132\\
23121223231322312132323122213233211112223323122121223332322112132121112231\\
21122233232332332323322233233211111221111123131222213122313122121322231211\\
22233222112132223121221313221312222131321111122111112332332223323233233232\\
33222112132211121231211223233322121221323322211112332312221323231213223132\\
32212132231312211121221111122332322323122233232232312221313221213121112122\\
33332221
\end{quote}

\section{Artifact inventory}

\medskip\noindent\textbf{\code{results/order94\_t12.template}.}\par\nobreak
Canonical order-$94$ certificate.  SHA-256:
\begin{center}
\hashbreak{a3dd3415956277d68fc0197f6f5e35c3}
{8f2ed7936c66f11f658b0c57688abec1}.
\end{center}

\medskip\noindent\textbf{\code{sources/rowley\_exoo\_order453.prototype}.}\par\nobreak
Source-recovered order-$453$ prototype.  SHA-256:
\begin{center}
\hashbreak{19c97e6279c184f6f462786cadda4b7c}
{9773d870a5b680a04eb1503ef384a2d0}.
\end{center}

\medskip\noindent\textbf{\code{results/r5\_5\_order42077.linear-coloring}.}\par\nobreak
Expanded $42076$-symbol five-colour distance word.  SHA-256:
\begin{center}
\hashbreak{a97d5dce8a927db2889ba220119f9d4}
{f3d1b88ee24d441f82a803a195ae8028d}.
\end{center}

\medskip\noindent\textbf{\code{sources/arxiv-2203.13476-source.tar.gz}.}\par\nobreak
Fixed source archive containing the ancillary XML.  SHA-256:
\begin{center}
\hashbreak{5041096a33d6898c116a35a102e521fb}
{14fe872fb18c20ff7d30822ed4c396b2}.
\end{center}

\begin{thebibliography}{9}
\raggedright

\bibitem{rowley-general}
F. Rowley,
\newblock A generalised linear Ramsey graph construction,
\newblock \emph{arXiv:1912.01164v3}, 2021.
\newblock See Definition 3.1 and Theorem 3.2 with proof, printed pp. 4--6.

\bibitem{rowley-sat}
F. Rowley,
\newblock Improved lower bounds for multicolour Ramsey numbers using
SAT-solvers,
\newblock \emph{arXiv:2203.13476v3}, 2022.
\newblock See Section 5, Table 1, and ancillary sheet
\code{Paper\_Sep\_2022}.

\bibitem{radziszowski}
S. P. Radziszowski,
\newblock Small Ramsey numbers,
\newblock \emph{Electronic Journal of Combinatorics, Dynamic Survey DS1},
revision 18, 24 April 2026.
\newblock DOI: \href{https://doi.org/10.37236/21}{10.37236/21}.

\bibitem{campos-pohoata}
M. Campos and C. Pohoata,
\newblock An update on multicolor Ramsey lower bounds,
\newblock \emph{arXiv:2601.15183v1}, 2026.

\end{thebibliography}

\end{document}
